\documentclass{article}
\usepackage{amsmath, amssymb, amsthm}
\usepackage{hyperref}
\usepackage{geometry}
\geometry{margin=1in}

\title{Erd\H{o}s Problem \#1204}
\date{Last updated: 2026-04-04}
\author{Source: erdosproblems.com}

\begin{document}
\maketitle

\noindent\textbf{Status:} OPEN \quad \textbf{No prize}

\noindent\textbf{Tags:} number theory

\medskip

\noindent\textbf{Problem Statement:}

We call a sequence of integers $0\leq a_1&#60;\cdots &#60;a_k$ admissible if it is missing at least one congruence class modulo every prime $p$. Let $A(k)=\min a_k$. Estimate $A(k)$ - in particular, is it true that\[A(k)\sim k\log k?\]Estimate\[B(k)=\min \frac{a_1+\cdots+a_k}{k}.\]




Erd\H{o}s \cite{Er80} attributes this problem to Elliott. The bounds are misstated in \cite{Er80}, but it is known that\[(1/2+o(1))k\log k \leq A(k)\leq (1+o(1))k\log k.\]The upper bound Erd\H{o}s attributes to Davenport, and follows simply by taking the $k$ smallest primes $&gt;k$. The lower order terms have been improved by Hensley and Richard \cite{HeRi73} (see Section 10 of \cite{Po14c} for more details). The lower bound is originally due to Elliott \cite{El65}, but was rediscovered by the Polymath project on bounded gaps between primes (see \cite{Po14c}). An upper bound of $\pi(x+y)\leq \pi(x)+(1+o(1))\pi(y)$ (see [855] ) together with the prime tuples conjecture would imply $A(k)\geq (1+o(1))k\log k$. It is trivial that $B(k)&lt;A(k)$. Taking the first $k$ primes which are $&gt;k$ implies\[B(k)\leq (1/2+o(1))k\log k,\]and since $a_j\geq A(j)$ we have $B(k) \geq \frac{1}{k}\sum_{1\leq j\leq k}A(j)$, whence if $A(k)\geq (c-o(1))k\log k$ we have $B(k) \geq (c/2-o(1))k\log k$, and so it is likely that\[B(k)\sim(1/2+o(1))k\log k.\]In \cite{Er80} Erd\H{o}s also asks about the behaviour of $a_k$ where $a_1&lt;a_2&lt;\cdots$ is the greedy admissible sequence (i.e. $a_1=0$ and $a_{i+1}$ is the smallest integer $&gt;a_{i-1}$ such that $a_1,\ldots,a_i$ is admissible).




References

[El65] Elliott, P. D. T. A., On sequences of integers . Quart. J. Math. Oxford Ser. (2) (1965), 35--45.

[Er80] Erd\H{o}s, Paul, A survey of problems in combinatorial number theory . Ann. Discrete Math. (1980), 89-115.

[HeRi73] Hensley, Douglas and Richards, Ian, On the incompatibility of two conjectures concerning primes . (1973), 123--127.

[Po14c] Polymath, D. H. J., Variants of the Selberg sieve, and bounded intervals containing many primes . Res. Math. Sci. (2014), Art. 12, 83.

\noindent\textbf{Related OEIS sequences:} A008407, A023193, A135311, possible


\bigskip
\noindent\small{Source: \url{https://www.erdosproblems.com/1204}}
\end{document}
